% Generated draft: requires mathematical and source review.
\documentclass[11pt]{article}
\usepackage[margin=25mm]{geometry}
\usepackage{fontspec}
\setmainfont{Latin Modern Roman}
\usepackage{amsmath,amssymb,mathtools,unicode-math}
\setmathfont{Latin Modern Math}
\usepackage{enumitem,needspace,xurl,hyperref}
\setlength{\parindent}{0pt}
\setlength{\parskip}{5pt}
\setlength{\emergencystretch}{3em}
\newcommand{\N}{\mathbb{N}}
\newcommand{\Z}{\mathbb{Z}}
\newcommand{\Q}{\mathbb{Q}}
\newcommand{\R}{\mathbb{R}}
\newcommand{\C}{\mathbb{C}}
\newcommand{\F}{\mathbb{F}}
\newcommand{\A}{\mathbb{A}}
\newcommand{\PP}{\mathbb{P}}
\newcommand{\E}{\mathbb{E}}
\newcommand{\eps}{\varepsilon}
\newcommand{\dd}{\,\mathrm{d}}
\begin{document}
\section*{UnsolvedMath ID: 701371}
\textbf{Source problem code:} catalogue entry
\subsection{Definitions and mathematical statement}
Let \(p\) be a prime. A group \(G\) has exponent \(p\) if \(G\neq\{1\}\) and \(g^p=1\) for every \(g\in G\). Let \(A\neq\{0\}\) be an abelian group, written additively. An action by automorphisms is a homomorphism \(\alpha:G\to\operatorname{Aut}(A)\). Such an action is fixed-point-free on the nonzero elements if
\[
\forall g\in G\setminus\{1\}\;\forall a\in A\setminus\{0\},\qquad \alpha(g)(a)\neq a.
\]
Under this interpretation of ``acting freely,'' the question proposed by V. D. Mazurov is: for every prime \(p\), every group \(G\) of exponent \(p\), every nontrivial abelian group \(A\), and every action \(\alpha\) satisfying the displayed condition, must \(G\) be cyclic, equivalently \(G\cong\Z/p\Z\)? No finiteness assumption is imposed on either group.

The supplied statement does not explicitly specify an action by automorphisms or say that freeness concerns only nonzero elements. These are unresolved interpretive qualifications of the formulation above. Every automorphism fixes \(0\), so freeness on the whole underlying set of \(A\) would be a different condition and would exclude every nontrivial acting group when the action is by automorphisms.
\subsection{Short English statement}
Under the fixed-point-free automorphism interpretation, must a group of prime exponent that acts on a nontrivial abelian group, with no nonidentity element fixing a nonzero element, be cyclic? The intended meaning of ``acting freely'' requires confirmation.
\subsection{Sources}
\begin{itemize}
\item Local UnsolvedMath registry, permanent ID \(701371\), exact source code \texttt{catalogue entry}, title \emph{Kourovka Notebook Problem 17.69}. The supplied record labels the problem open; this status is not independently verified.
\item \emph{The Kourovka Notebook (Unsolved Problems in Group Theory)}, Problem 17.69, proposed by V. D. Mazurov; 17th issue (2010). The supplied background also identifies the 21st issue edition (2026), arXiv identifier \texttt{1401.0300v48}.
\item The supplied collection description identifies E. I. Khukhro and V. D. Mazurov as editors of the notebook collection covering Issues 1--21 (1965--2026). These bibliographic details are supplied attribution and have not been independently checked.
\end{itemize}
\end{document}
